\documentclass[12pt]{article}
\usepackage[utf8]{inputenc}
\usepackage[T1]{fontenc}
\usepackage[margin=1in]{geometry}
\usepackage{setspace}
\usepackage{xcolor}
\usepackage[colorlinks=true, urlcolor=blue!60!black, linkcolor=blue!60!black]{hyperref}
\setlength{\parindent}{2em}
\setlength{\parskip}{1.25em}
\renewcommand{\baselinestretch}{1.0}
\begin{document}
\begin{center}
{\LARGE \textit{Allen v.\ Wright}, 468 U.S.\ 737 (1984)\par}
\vspace{0.8em}
{\large Case Brief --- Four-Page Edition\par}
\end{center}
\vspace{1em}

\section{Memory Jogger}

Parents of Black public-school children lacked Article~III standing to sue the IRS over lax enforcement of the rule denying tax exemptions to racially discriminatory private schools: abstract stigmatic harm is not cognizable without personal denial of equal treatment, and any impairment of desegregation is not ``fairly traceable'' to the IRS (\href{https://supreme.justia.com/cases/federal/us/468/737/}{Justia}).

\section{Detailed Case Facts}

After \textit{Brown v.\ Board of Education}, private ``segregation academies'' opened or expanded in communities undergoing desegregation. In 1970--1971 the IRS adopted a policy denying tax-exempt status under \S~501(c)(3) to racially discriminatory private schools, and Rev.\ Proc.\ 75-50 required schools to adopt, publicize, and certify a nondiscriminatory policy. In 1976, Inez Wright and other parents of Black children in desegregating districts in seven States filed a nationwide class action against the Treasury Secretary and the IRS Commissioner; W.\ Wayne Allen, head of a named Memphis school system, intervened. Plaintiffs alleged the guidelines let discriminatory schools obtain exemptions by certifying a paper policy. They never alleged their children had applied or would apply to any private school. They claimed (1) direct harm and stigma from federal aid to discriminatory schools and (2) impairment of their children's opportunity for a desegregated public education, and sought an injunction forcing stricter IRS rules. Meanwhile, Congress blocked tougher IRS procedures through appropriations riders for fiscal years 1980--1982.

\section{Procedural History}

The district court dismissed for lack of standing. A divided D.C.\ Circuit reversed, finding standing in the ``denigration'' injury to Black parents and children and enjoining the IRS. The Supreme Court granted the petitions of the Government and of intervenor Allen and reversed.

\section{Judicial Votes}

\textbf{5--3}, reversing and vacating the injunction (July 3, 1984).
\begin{itemize}
\item \textbf{Majority:} Justice O'Connor, joined by Chief Justice Burger and Justices White, Powell, and Rehnquist.
\item \textbf{Dissents:} Justice Brennan; Justice Stevens, joined by Justice Blackmun.
\item \textbf{Recused:} Justice Marshall took no part.
\end{itemize}

\section{Holding}

Plaintiffs lack Article~III standing. Standing requires a personal injury, fairly traceable to the defendant's allegedly unlawful conduct, and likely to be redressed by the relief requested. A stigmatic injury from racial discrimination confers standing only on those \textit{personally denied equal treatment}; and the diminished-desegregation injury, which depends on independent choices of private schools and white parents, is too attenuated to be traceable to the IRS---especially where relief would have courts restructure Executive enforcement.

\section{Analysis of Opinions}

\textbf{Majority (O'Connor).} Standing rests on separation of powers. An asserted right to have the Government obey the law does not suffice, and precedent limits stigmatic-injury standing to persons ``personally denied equal treatment''; otherwise federal courts would become forums for ``concerned bystanders.'' The children's diminished ability to receive a desegregated education was ``beyond any doubt'' a serious, cognizable injury, but under \textit{Simon v.\ Eastern Kentucky Welfare Rights Org.} the causal chain was fatally attenuated: it was uncertain how many discriminatory schools held exemptions, and speculative whether losing them would change school policies or send white children back to public schools. Recognizing standing would make courts ``continuing monitors'' of the Executive. Footnote~19 separated traceability from redressability, a distinction later canonized in \textit{Lujan v.\ Defenders of Wildlife}.

\textbf{Brennan (dissenting).} The majority used ``standing to slam the courthouse door.'' On a motion to dismiss, the allegations must be accepted as true, and the complaint named at least 32 discriminatory schools, 18 in Memphis alone. \textit{Norwood v.\ Harrison} and \textit{Coit v.\ Green} established that aid to segregated academies impedes desegregation.

\textbf{Stevens, joined by Blackmun (dissenting).} Injury was undisputed; causation followed from ``elementary economics,'' since tax exemptions function as ``cash grants'' and revoking them raises the cost of segregated schooling. The majority conflated standing with merits questions about executive discretion, which \textit{Marbury} makes a judicial function.

\section{Examples: Future Applications}

\begin{enumerate}
\item \textbf{Same Side: Enforcement Priorities.} Two States sue to make federal officers arrest more noncitizens. In \textit{United States v.\ Texas} (2023), the Court, paralleling \textit{Allen}, held that plaintiffs who are not the object of enforcement lack standing to challenge arrest priorities.
\item \textbf{Same Side: Speculative Third Parties.} Environmentalists challenge a rule on foreign-project consultation, hoping to see endangered species ``some day.'' \textit{Lujan} extended \textit{Allen} to deny standing.
\item \textbf{Opposite Side: Personal Rejection.} A Black student is personally rejected by a tax-exempt academy because of her race. She was personally denied equal treatment and has standing.
\item \textbf{Opposite Side: Unequal Barrier.} Contractors challenge a minority set-aside without proving they would have won a bid. Under \textit{Northeastern Florida} (1993), being subjected to an unequal barrier is itself a traceable injury.
\item \textbf{Fence-Sitter: Market Predictions.} Fuel producers challenge EPA's approval of California's electric-vehicle mandate. The injury flows through automakers and consumers, yet \textit{Diamond Alternative Energy v.\ EPA} (2025) found commonsense economic predictability enough---echoing Justice Stevens's logic.
\end{enumerate}

\section{Critique}

\textbf{Merits in disguise.} Scholars such as Gene Nichol argue the Court manipulated causation to duck a sensitive merits dispute, imposing a summary-judgment burden at the pleading stage, and note the asymmetry with later cases that readily credited market incentives.

\textbf{Structural asymmetry and originalism.} Progressive scholars such as Cass Sunstein argue \textit{Allen} lets regulated entities challenge enforcement while locking out regulatory beneficiaries when agencies underenforce. Then-Judge Scalia celebrated its separation-of-powers thesis, but newer originalist work calls the injury-traceability trilogy ``faux originalism.''

\section{Key Quotations}

\begin{quote}
``More important, the law of Art.\ III standing is built on a single basic idea --- the idea of separation of powers.''
\end{quote}
\noindent --- Justice O'Connor, for the Court, 468 U.S.\ at 752.

\begin{quote}
``If the abstract stigmatic injury were cognizable, standing would extend nationwide to all members of the particular racial groups against which the Government was alleged to be discriminating by its grant of a tax exemption to a racially discriminatory school, regardless of the location of that school. All such persons could claim the same sort of abstract stigmatic injury respondents assert in their first claim of injury. A black person in Hawaii could challenge the grant of a tax exemption to a racially discriminatory school in Maine.''
\end{quote}
\noindent --- Justice O'Connor, for the Court, 468 U.S.\ at 755--756.

\begin{quote}
``The Constitution, after all, assigns to the Executive Branch, and not to the Judicial Branch, the duty to `take Care that the Laws be faithfully executed.'\,''
\end{quote}
\noindent --- Justice O'Connor, for the Court, 468 U.S.\ at 761.

\end{document}
